% Conversão fiel da participação como cliente Hoortech; origem imutável RELATORIO_ATUAL.odt.
\hypertarget{__RefHeading___Toc4470_311817663332}{}

% SOURCE text-h-0035
\refstepcounter{subsection}
\subsection*{6.2.6 Sprint 4}
\label{cliente-hoortech:text-h-0035}
\addcontentsline{toc}{subsection}{6.2.6 Sprint 4}

% SOURCE text-p-0742
Na Sprint 4, tive a oportunidade de acompanhar a entrega final do projeto desenvolvido pelo aluno ao longo de um semestre inteiro. A entrega foi apresentada a nós, stakeholder, antes da grande apresentação final para todos os grupos da AGES. Fiquei profundamente surpreendido com o resultado alcançado pelo time. O projeto estava praticamente impecável, demonstrando um grande salto em termos de funcionalidades e qualidade desde as primeiras Sprints.\par

% SOURCE text-p-0743
Nessa última Sprint, a equipe implementou o alfabeto completo no modelo, permitindo que ele reconhecesse e traduzisse os sinais correspondentes com precisão. Além disso, alguns bugs pendentes de Sprints anteriores foram resolvidos, contribuindo para um refinamento geral do sistema. O progresso técnico foi notável, mas uma pendência permaneceu: o grupo enfrentou dificuldades em integrar o back-end ao front-end no servidor da AWS devido a um problema relacionado ao uso do WebSocket. Apesar de se comprometerem a resolver a questão antes da apresentação final, o desafio se mostrou mais complexo do que o esperado, e a integração não foi concluída a tempo. A apresentação final foi realizada com o front-end rodando no servidor local, porém se comunicando com o modelo em uma EC2 na AWS.\par

% SOURCE text-p-0744
Essa limitação, no entanto, não tirou o brilho da apresentação final realizada pelo grupo. O time conduziu uma demonstração impecável, evidenciando todo o potencial e a relevância do projeto. A clareza na apresentação e a qualidade do produto foram destacadas, mostrando o comprometimento e o profissionalismo do grupo ao longo do desenvolvimento.\par

% SOURCE text-p-0745
Após a apresentação, fui convidado pelo professor coordenador da AGES, Marcelo Yamaguti, a comentar sobre o projeto. Ressaltei a complexidade e os desafios enfrentados pela equipe ao longo das Sprints, bem como os limitadores que tornaram o projeto ainda mais exigente. Destacar esses pontos foi importante para evidenciar o mérito do grupo em superar tantas barreiras com criatividade e dedicação.\par

% SOURCE text-p-0746
Ao final da apresentação, fomos agraciados com a honraria de “Projeto Destaque” da sexta-feira na edição 2024/2 da AGES. Foi um momento de enorme felicidade e realização, não só para mim enquanto cliente, mas especialmente para o time. Essa conquista é, acima de tudo, o reflexo do trabalho árduo, da confiança mútua e do compromisso do grupo com a excelência. Ver o empenho e o talento deles serem reconhecidos foi extremamente gratificante, e não tenho dúvidas de que o sucesso alcançado nesse projeto é um prenúncio de muitas outras conquistas no futuro.\par

% FIGURE draw-image-0007; source text-p-0749; image provenance pending human identification.
\memorialfigure{imagens/original/asset-0016.jpg}{\textbf{Figura 33 –} Entrega Final - CLIENTE AGES}{cliente-hoortech:draw-image-0007}{Fonte: pendente de identificação}{0.92}

